%% Chương 4: Các thuật toán so sánh
\clearpage
\phantomsection
\section*{CHƯƠNG 4. CÁC THUẬT TOÁN SO SÁNH}
\addcontentsline{toc}{section}{\numberline{}CHƯƠNG 4. CÁC THUẬT TOÁN SO SÁNH}
\label{ch:baseline}
\setcounter{section}{4}
\setcounter{subsection}{0}
\setcounter{figure}{0}
\setcounter{table}{0}
\setcounter{equation}{0}

\subsection{Tổng quan}

Để đánh giá khách quan thuật toán hybrid đề xuất, ba thuật toán mã hoá cổ điển
được cài đặt và tích hợp vào cùng pipeline với cùng interface IEncryption.
Việc so sánh trên cùng dataset, cùng chiến lược (S1/S2), và cùng bộ chỉ số
đảm bảo tính công bằng (apples-to-apples comparison).

\textbf{Nguyên tắc thiết kế baseline:}
Ba thuật toán này được cố ý thiết kế với \textbf{keystream cố định}
(fixed keystream cho mọi NALU) để làm nổi bật điểm yếu của thiết kế không có
per-NALU variation, qua đó chứng minh giá trị của tính năng per-NALU keystream
trong thuật toán hybrid.

Các điểm yếu được giới thiệu dưới đây không phải là bản chất yếu kém của AES/DES/RC4
như thuật toán mật mã --- mà là hệ quả của \emph{chế độ sử dụng cụ thể} trong
ngữ cảnh mã hoá chọn lọc với fixed keystream. AES được chuẩn hoá FIPS và vẫn
là thuật toán cực kỳ an toàn khi sử dụng đúng cách (với IV/counter ngẫu nhiên).

\subsection{AES-128-CTR (Advanced Encryption Standard --- Counter Mode)}

\subsubsection{Cài đặt}

AES-128 với chế độ CTR, counter = 0 cố định cho mọi NALU. Khóa K được
pad/truncate về đúng 16 bytes (128-bit) bằng cách:
\begin{itemize}
  \item Nếu $|K| < 16$: pad zero bên phải.
  \item Nếu $|K| > 16$: truncate hoặc XOR fold.
\end{itemize}

Với 25 bytes DC, quá trình CTR sinh keystream:
\begin{itemize}
  \item Block 1: $ks[0..15] = \text{AES}_K(\texttt{counter=0} \| \texttt{nonce=0})$ (16 bytes)
  \item Block 2: $ks[16..31] = \text{AES}_K(\texttt{counter=1} \| \texttt{nonce=0})$ (16 bytes, lấy 9 bytes đầu)
  \item Tổng: 2 lần gọi AES per NALU.
\end{itemize}

Ciphertext: $C[i] = \text{DC}[i] \oplus ks[i]$ cho $i = 0, \ldots, 24$.

\subsubsection{Phân tích điểm yếu: Keystream reuse}

Do nonce = 0 và counter bắt đầu từ 0 cố định cho mọi NALU, keystream là hằng số:
\begin{equation}
ks = \text{AES}_K(0) \| \text{AES}_K(1) \quad \text{(giống nhau cho mọi NALU)}
\end{equation}

\textbf{Hậu quả 1 --- Known-plaintext attack:}
Nếu kẻ tấn công có được plaintext $P_1$ (DC bytes của một NALU qua phân tích
thống kê hoặc kênh bên), hắn tính được:
\begin{equation}
ks = C_1 \oplus P_1
\end{equation}
Sau đó giải mã mọi NALU khác: $P_j = C_j \oplus ks$.

\textbf{Hậu quả 2 --- IND-CPA violation:}
Hai NALUs có cùng DC bytes ($P_1 = P_2$) cho ciphertext giống hệt ($C_1 = C_2$).
Kẻ tấn công phân biệt được ciphertext của cùng plaintext, vi phạm IND-CPA.

\textbf{Hậu quả 3 --- XOR of ciphertexts:}
\begin{equation}
C_j \oplus C_k = (P_j \oplus ks) \oplus (P_k \oplus ks) = P_j \oplus P_k
\end{equation}
Kẻ tấn công biết XOR của hai plaintext không cần biết khóa.

\subsubsection{Phân tích NPCR/UACI cho AES-CTR fixed}

Phép đo NPCR của Wu \cite{wu_npcr} yêu cầu so sánh $C_1 = \text{Enc}(P)$
và $C_2 = \text{Enc}(P')$ với $P'[0] = P[0] \oplus 1$.

Với AES-CTR fixed ($ks$ hằng số):
\begin{align}
C_1[i] &= P[i] \oplus ks[i] \\
C_2[i] &= P'[i] \oplus ks[i]
\end{align}

$C_1[i] \neq C_2[i]$ khi và chỉ khi $P[i] \neq P'[i]$, tức là khi $i = 0$ (byte bị flip).

\begin{equation}
\text{NPCR}_{\text{AES-fixed}} = \frac{1}{25} \times 100\% = 4\%
\end{equation}

Giá trị 4\% này xác định bởi cấu trúc toán học (1/25 bytes thay đổi), không có
ý nghĩa đánh giá bảo mật. Do đó NPCR/UACI được ghi là ``N/A'' cho AES trong báo cáo.

\subsubsection{Lý do giữ lại AES trong so sánh}

Mặc dù có điểm yếu về keystream reuse, AES-128-CTR được giữ lại vì:
\begin{enumerate}
  \item AES là chuẩn mật mã được chấp thuận rộng rãi (FIPS 197, NIST), tốc độ cao
        nhờ AES-NI hardware instruction trên CPU x86-64 hiện đại.
  \item Kết quả so sánh chứng minh rằng \emph{chế độ hoạt động và thiết kế IV}
        quan trọng hơn độ mạnh của block cipher. Dùng AES-CTR đúng cách
        (IV ngẫu nhiên mỗi NALU) sẽ không có điểm yếu này.
  \item AES đại diện cho nhóm ``thuật toán tiêu chuẩn nhưng thiếu per-NALU variation''.
\end{enumerate}

\subsection{DES-OFB (Data Encryption Standard --- Output Feedback Mode)}

\subsubsection{Cài đặt}

DES \cite{fips_des} với chế độ OFB, IV = 0 cố định. Khóa K được rút gọn về
8 bytes (64-bit raw key, 56-bit hiệu dụng sau khi DES bỏ parity bits ở vị trí
$i \equiv 0 \pmod{8}$).

Sinh keystream OFB:
\begin{align}
O_0 &= IV = 0 \\
O_j &= \text{DES}_K(O_{j-1}), \quad j = 1, 2, \ldots \\
C[8(j-1)..8j-1] &= \text{DC}[8(j-1)..8j-1] \oplus O_j
\end{align}

DES block size = 8 bytes, cần 4 block cho 25 bytes ($\lceil 25/8 \rceil = 4$,
lấy 25 bytes đầu của 32 bytes keystream).

\subsubsection{Điểm yếu bảo mật}

\textbf{Điểm yếu 1 --- Keystream cố định (giống AES-CTR):}
IV = 0 và K cố định → keystream $O_1, O_2, O_3, O_4$ hằng số → tất cả hệ quả
của AES-CTR fixed cũng áp dụng cho DES-OFB.

\textbf{Điểm yếu 2 --- Khóa 56-bit đã lỗi thời:}
Năm 1998, EFF DES Cracker (máy phần cứng chuyên dụng \$250.000) phá vỡ DES
trong chưa đầy 23 giờ bằng exhaustive key search. Với GPU hiện đại có thể
thử $\sim10^{12}$ khóa/giây, brute-force 56-bit ($\approx 7.2 \times 10^{16}$
khóa) cần khoảng 20 ngày.

\textbf{Điểm yếu 3 --- Birthday attack trong OFB:}
Do OFB là keystream generator tuyến tính trên $\mathbb{Z}_2^{64}$, chu kỳ lý thuyết
là $2^{64}$. Tuy nhiên, birthday paradox cho thấy sau $\sim 2^{32}$ block ($2^{35}$ bytes),
có xác suất 50\% keystream cycle lặp lại --- gây keystream reuse cấp block,
nguy hiểm hơn keystream reuse cấp NALU.

\textbf{Điểm yếu 4 --- Software DES chậm:}
DES không có hardware acceleration trên x86-64 (khác với AES có AES-NI).
Cài đặt phần mềm phải dùng S-box lookup table (8 S-boxes $\times$ 6-bit→4-bit),
dẫn đến throughput thấp hơn AES khoảng 8 lần.

\subsubsection{Anomaly hiệu năng: DES S1 Decrypt}

Trong thực nghiệm (phân tích chi tiết ở Chương 6), DES S1 decrypt pipeline time
đo được 8.686 s --- bất thường so với encrypt 1.519 s. Tuy nhiên, DC decrypt time
thực tế = 478 ms (nhất quán với encrypt 447 ms). Nguyên nhân: disk I/O page cache
miss do pipeline decrypt đọc file mới ($\neq$ file vừa encrypt được cache). Đây là
artifact hệ thống, không phải lỗi thuật toán.

\subsection{RC4 (Rivest Cipher 4 --- Per-NALU Reinit)}

\subsubsection{Cài đặt}

RC4 được khởi tạo lại hoàn toàn từ đầu với mỗi NALU sử dụng cùng khóa K.
Mỗi NALU: chạy KSA (256 byte, 256 hoán vị), sau đó lấy 25 bytes đầu của PRGA.

Không giống AES và DES, RC4 là stream cipher thuần túy --- không có khái niệm
``block'' hay ``IV''. Keystream hoàn toàn xác định bởi khóa K (với RC4 standard).

\subsubsection{Phân tích tốc độ}

RC4 nhanh nhất trong 4 thuật toán do cấu trúc cực kỳ đơn giản:
\begin{itemize}
  \item KSA: 256 swap operations $\times$ 1 NALU = 256 operations.
  \item PRGA: 25 iterations (3 operations/iteration) = 75 operations.
  \item Tổng: $\approx$ 331 operations/NALU, so với AES: $\sim$2000 operations/NALU.
\end{itemize}

Kết quả thực nghiệm: RC4 đạt 769 ns/NALU S1, nhanh hơn hybrid $38\times$ và AES $3.6\times$.

\subsubsection{Điểm yếu bảo mật}

\textbf{Điểm yếu 1 --- Keystream cố định:}
Vì mỗi NALU đều khởi tạo RC4 từ đầu với cùng K, 25 bytes PRGA output đầu tiên
\textbf{giống hệt nhau cho mọi NALU}. Đây là keystream reuse nghiêm trọng nhất
trong 3 baseline (vì RC4 không có IV hay counter).

\textbf{Điểm yếu 2 --- Bias thống kê trong output đầu:}
Fluhrer-Mantin-Shamir attack \cite{fms_rc4} (2001) chứng minh rằng byte 0 của
RC4 PRGA output có thiên lệch thống kê: byte 0 có xác suất cao hơn là \texttt{K[1] + K[2]}.
Điều này làm 25 bytes đầu đặc biệt dễ bị tấn công phân tích thống kê.

Phiên bản WPA (Wi-Fi Protected Access) đã bị tấn công thành công bằng cách khai
thác thiên lệch này, dẫn đến việc RC4 bị cấm trong TLS 1.3 và nhiều giao thức hiện đại.

\textbf{Điểm yếu 3 --- Thiếu diffusion:}
Tương tự AES-CTR và DES-OFB, RC4 không có cơ chế khuếch tán giữa các bytes.
1 bit thay đổi plaintext → 1 bit thay đổi ciphertext → NPCR $\approx 4\%$.

\subsection{So sánh định tính 4 thuật toán}

\begin{table}[ht]
\centering
\caption{So sánh định tính 4 thuật toán về các tiêu chí thiết kế}
\label{tab:qualitative}
\begin{tabular}{lccccc}
\toprule
Tiêu chí & Hybrid & AES-CTR & DES-OFB & RC4 \\
\midrule
Keystream per-NALU   & \checkmark & $\times$ & $\times$ & $\times$ \\
Diffusion giữa bytes & \checkmark & $\times$ & $\times$ & $\times$ \\
NPCR $> 99.6\%$      & \checkmark & N/A & N/A & N/A \\
Entropy $> 7.9$      & \checkmark & \checkmark & \checkmark & \checkmark \\
Keyspace $> 2^{56}$  & \checkmark ($2^{58}$) & \checkmark ($2^{128}$) & $\times$ ($2^{56}$) & \checkmark ($2^{128}$) \\
Tốc độ DC S1 (ns/NALU) & 29.245 & 2.753 & 23.500 & 769 \\
Chuẩn hoá            & $\times$ & FIPS 197 & FIPS 46-3 & $\times$ \\
Hardware acceleration & $\times$ & AES-NI & $\times$ & $\times$ \\
\bottomrule
\end{tabular}
\end{table}

\textbf{Nhận xét tổng hợp:}
\begin{itemize}
  \item Hybrid là thuật toán duy nhất có đầy đủ per-NALU keystream variation
        và intra-NALU diffusion --- hai tính chất cốt lõi để đạt NPCR $> 99.6\%$.
  \item AES là lựa chọn tốt nhất trong 3 baseline: chuẩn hoá, tốc độ cao, keyspace lớn.
        Thiếu điểm duy nhất là per-NALU keystream và diffusion.
  \item RC4 nhanh nhất nhưng có thêm điểm yếu thống kê (FMS attack), không được
        khuyến nghị cho ứng dụng mới.
  \item DES không khuyến nghị: keyspace $2^{56}$ đã bị phá vỡ, chậm (thiếu hardware acceleration).
\end{itemize}

\subsection{Phân tích toán học: Tại sao NPCR/UACI không áp dụng cho AES/DES/RC4}

\subsubsection{Chứng minh hình thức}

Gọi $P$ là DC bytes gốc và $P' = P \oplus e_0$ với $e_0 = (1, 0, \ldots, 0)$ (flip bit 0 byte 0).
Gọi $ks$ là keystream cố định (giống nhau cho mọi NALU và mọi plaintext).

\begin{align}
C_1[i] &= P[i] \oplus ks[i] \label{eq:c1} \\
C_2[i] &= P'[i] \oplus ks[i] = (P[i] \oplus e_0[i]) \oplus ks[i] \label{eq:c2}
\end{align}

Trong đó $e_0[i] = 1$ nếu $i = 0$, $e_0[i] = 0$ nếu $i \neq 0$.

Hiệu:
\begin{equation}
C_1[i] \oplus C_2[i] = e_0[i] = \begin{cases} 1 & i = 0 \\ 0 & i \neq 0 \end{cases}
\end{equation}

Do đó $D[i] = 1$ chỉ khi $i = 0$, và $D[i] = 0$ với mọi $i \neq 0$.

\begin{equation}
\text{NPCR}_{\text{fixed}} = \frac{\sum_{i=0}^{24} D[i]}{25} \times 100\%
= \frac{1}{25} \times 100\% = 4\%
\end{equation}

\begin{equation}
|C_1[0] - C_2[0]| = |(P[0] \oplus ks[0]) - (P'[0] \oplus ks[0])|
\end{equation}

Vì $P'[0] = P[0] \oplus 1$: nếu bit 0 của $P[0] \oplus ks[0]$ là 0 thì $|C_1[0] - C_2[0]| = 1$;
nếu là 1 thì $= -1$ về giá trị unsigned, tức là $= 255$ (mod 256).

\begin{equation}
\text{UACI}_{\text{fixed}} = \frac{|C_1[0] - C_2[0]|}{25 \times 255} \times 100\%
= \frac{1 \text{ hoặc } 255}{25 \times 255} \times 100\% \in \{0.016\%, 4\%\}
\end{equation}

Cả hai giá trị này xa ngưỡng $33.46\%$ và hoàn toàn xác định bởi 1 byte.
Chúng không phản ánh bất kỳ tính chất bảo mật nào, do đó NPCR/UACI được ghi
là ``N/A'' (không áp dụng) trong tất cả bảng kết quả.

\subsubsection{Ngụ ý cho thiết kế}

Kết quả trên chứng minh:
\begin{enumerate}
  \item NPCR $> 99.6\%$ \textbf{không thể đạt được} với bất kỳ thuật toán XOR nào
        có keystream cố định, bất kể độ mạnh của block/stream cipher.

  \item Để đạt NPCR $> 99.6\%$, \textbf{bắt buộc} phải có hoặc:
        (a) diffusion feedback (chained diffusion như hybrid), hoặc
        (b) keystream phụ thuộc vào ciphertext trước (như CBC mode), hoặc
        (c) cả hai.

  \item Chained diffusion của hybrid giải quyết (a), cho phép đạt NPCR $> 99.6\%$
        ngay cả khi chỉ mã hoá 25 bytes.
\end{enumerate}

\subsection{Tổng kết Chương 4}

Ba thuật toán baseline AES-128-CTR, DES-OFB, và RC4 trong cấu hình fixed-keystream
đều có chung một điểm yếu cốt lõi: \emph{không có per-NALU keystream variation và
không có intra-NALU diffusion}. Kết quả là NPCR $\approx 4\%$, vi phạm IND-CPA,
và dễ bị tấn công known-plaintext.

Các điểm yếu này không phải nhược điểm cố hữu của AES/DES/RC4 như thuật toán mật mã
--- mà là hệ quả của thiết kế keystream cố định trong ngữ cảnh mã hoá chọn lọc NALU.
Thuật toán hybrid đề xuất giải quyết các điểm yếu này bằng PLCM per-NALU keystream
và chained diffusion feedback, đạt NPCR $> 99.6\%$ (S2) và Entropy $\approx 8.0$.
